\chapter{Estado del arte}
\label{chap:estadoarte}

\lettrine{L}{a} gestión digital de instalaciones deportivas es un mercado en
rápido crecimiento, impulsado por la popularización del pádel y la demanda de
servicios \textit{on-demand}. En este capítulo se analizan las soluciones
existentes, se evalúan las tecnologías candidatas para el desarrollo del sistema
y se justifica la elección tecnológica adoptada.

\section{Soluciones existentes}
\label{sec:soluciones_existentes}

\subsection{Playtomic}

Playtomic es la plataforma líder en España para la reserva de pistas de pádel y
tenis. Permite a los usuarios buscar pistas disponibles en su zona, realizar
reservas en tiempo real y gestionar sus partidos. Desde el punto de vista del
club, ofrece un panel de administración para gestionar disponibilidad, precios y
bonos.

Sus principales fortalezas son la amplia red de clubes integrados y la experiencia
de usuario en la aplicación móvil. Sin embargo, presenta limitaciones relevantes:
es una solución SaaS cerrada sin posibilidad de personalización, cobra una
comisión por reserva que puede resultar elevada para clubes pequeños, y no ofrece
funcionalidades avanzadas de gestión de torneos integradas con el sistema de
reservas.

\subsection{Smashapp}

Smashapp es una plataforma orientada a la gestión de ligas y torneos de pádel.
Su punto fuerte es el módulo de competición, que permite crear torneos, gestionar
inscripciones y publicar resultados. No obstante, sus capacidades de gestión de
reservas son limitadas y está orientado a competiciones federadas más que a la
gestión cotidiana de un club.

\subsection{Sistemas propios de clubes}

Muchos clubes han desarrollado o contratado sistemas a medida, a menudo
implementados como simples formularios web conectados a una base de datos o
incluso como hojas de cálculo compartidas. Estos sistemas presentan problemas
de mantenibilidad, falta de integración entre módulos y dependencia de un
proveedor concreto.

\subsection{Análisis comparativo}

La tabla~\ref{tab:comparativa_soluciones} resume las funcionalidades principales
de las soluciones analizadas.

\begin{table}[ht]
\centering
\resizebox{\textwidth}{!}{%
\begin{tabular}{lcccc}
\hline
\textbf{Funcionalidad} & \textbf{Playtomic} & \textbf{Smashapp} & \textbf{Sist.\ propios} & \textbf{PadelCenter} \\
\hline
Reserva de pistas        & Sí       & Limitado  & Variable & Sí \\
Gestión de torneos       & Básico   & Sí        & Raro     & Sí \\
Emparejamientos auto.    & No       & Sí        & No       & Sí \\
Personalizable           & No       & Limitado  & Sí       & Sí \\
Código abierto           & No       & No        & Variable & Sí \\
Sin comisión por reserva & No       & No        & Sí       & Sí \\
Arquitectura moderna     & Descono. & Descono.  & No       & Sí \\
\hline
\end{tabular}}
\caption{Comparativa de soluciones existentes.}
\label{tab:comparativa_soluciones}
\end{table}

\section{Tecnologías evaluadas}
\label{sec:tecnologias_evaluadas}

\subsection{Frameworks de backend}

Se evaluaron tres frameworks Java para el backend:

\paragraph{Spring Boot~\cite{springboot}} Es el framework más maduro y ampliamente adoptado del
ecosistema Java. Ofrece un ecosistema de módulos muy completo (Spring Security,
Spring Data JPA, Spring MVC), una comunidad enorme y excelente integración con
herramientas de generación de código a partir de OpenAPI. Su modelo de
programación es familiar para cualquier desarrollador Java.

\paragraph{Quarkus} Framework diseñado para entornos cloud-native, con tiempos
de arranque muy reducidos y menor huella de memoria. Aunque es una opción muy
interesante para microservicios, su ecosistema es menos maduro que Spring Boot
y la curva de aprendizaje es mayor para quienes provienen del mundo Java EE.

\paragraph{Micronaut} Similar a Quarkus en su orientación cloud-native, con
inyección de dependencias en tiempo de compilación. Sus ventajas en rendimiento
son similares a las de Quarkus, pero su adopción en el mercado es significativamente
menor, lo que implica menor disponibilidad de recursos y librerías.

La elección recayó sobre \textbf{Spring Boot} por su madurez, la calidad del
plugin OpenAPI Generator para Maven y la cobertura de todas las necesidades
del proyecto sin compromisos.

\subsection{Frameworks de frontend}

Se evaluaron tres frameworks JavaScript/TypeScript para el frontend:

\paragraph{Next.js~\cite{nextjs}} Framework React con renderizado híbrido (SSR, SSG y RSC),
App Router en su versión 14, y un ecosistema muy activo. Su modelo de enrutamiento
basado en el sistema de ficheros es intuitivo y facilita la organización del
código. La integración con React Query y la generación de código con Orval es
directa.

\paragraph{Nuxt} Equivalente de Next.js para el ecosistema Vue.js. Ofrece
características similares pero la comunidad de Vue es más pequeña y el autor
tiene mayor experiencia con React.

\paragraph{Angular} Framework completo de Google, con un enfoque más opinionado
y una mayor cantidad de boilerplate. Su curva de aprendizaje es más pronunciada
y su arquitectura, aunque sólida, resulta excesiva para un proyecto de este
tamaño.

La elección recayó sobre \textbf{Next.js} por su flexibilidad, la calidad del
ecosistema React y la compatibilidad con las herramientas de generación de código
elegidas.

\section{Arquitectura hexagonal vs. arquitectura en capas}
\label{sec:arquitectura_comparativa}

La arquitectura en capas tradicional (\textit{Presentation → Service → Repository})
es el patrón más extendido en aplicaciones Spring Boot. Sin embargo, presenta
un problema fundamental: las capas superiores acaban dependiendo de detalles de
infraestructura (JPA, SQL, HTTP) que se filtran a través de las capas, acoplando
el dominio de negocio a tecnologías concretas.

La arquitectura hexagonal, también conocida como \textit{Ports \& Adapters} y
propuesta por Alistair Cockburn~\cite{cockburn2005hexagonal}, resuelve este
problema invirtiendo las dependencias: el núcleo de la aplicación (dominio y
casos de uso) define interfaces (\textit{puertos}) que la infraestructura
implementa (\textit{adaptadores}). Esto garantiza que el dominio nunca depende
de JPA, Spring o cualquier otra tecnología concreta. Esta misma idea de
inversión de dependencias hacia el dominio es el principio central de la
\textit{Clean Architecture} de Robert C.~Martin~\cite{martin2017cleanarchitecture},
que PadelCenter sigue de forma directa al organizar el código en las capas
\texttt{domain}, \texttt{application} e \texttt{infrastructure}. El propio
patrón de mapeo en dos capas descrito en la sección~\ref{sec:implementacion_backend}
recoge ideas clásicas de mapeo objeto-relacional descritas por
Fowler~\cite{fowler2002peaa}, adaptadas para mantener las entidades JPA
completamente fuera del dominio.

Las ventajas concretas para PadelCenter son:
\begin{itemize}
  \item Posibilidad de testar los casos de uso sin levantar el contexto Spring.
  \item Facilidad para sustituir la base de datos o el mecanismo de autenticación
        sin modificar el dominio.
  \item Reglas de arquitectura verificables automáticamente con ArchUnit.
\end{itemize}

La desventaja principal es el mayor número de clases e interfaces necesarias
respecto a una arquitectura en capas plana. Sin embargo, en un proyecto con
ambición de mantenibilidad a largo plazo, esta inversión es justificada.
